\section{R panel software: compare the target before comparing the estimator}
\label{sec:exec24-panels-r}

A branch rollout can ask about a particular adoption cohort, a calendar
quarter, the first year after adoption, or the average benefit across all
exposed branch-periods. These are different targets even before nuisance
models or standard errors enter. The R comparison therefore connects the
selected \texttt{did} implementation to the already inspected
\texttt{DRDID}, \texttt{did2s}, \texttt{fixest}, \texttt{HonestDiD},
\texttt{synthdid} and \texttt{gsynth}/\texttt{fect} paths. These are concrete
alternatives with different assumptions and interfaces, not interchangeable
labels for an automatic causal answer.
\citep{exec24_panelr_contract}

The selected \texttt{did} revision is dated 28 August 2026 and declares
version 2.5.1.902. Its previously retained \texttt{att\_gt} file authenticates
to the same complete cutoff-pinned tree. We additionally inspect aggregation,
preprocessing, cell estimation and multiplier-bootstrap delegates. The
revision's declared GPL-3 license, R requirement and dependencies are
observations about these bytes, not proof of an installed distribution,
commercial support or qualification of every imported estimator.
\citep{exec24_panelr_did}

\subsection{Four aggregate answers from the same cohort-time effects}
Write $\tau_{gt}$ for the effect in cohort $g$ at calendar time $t$ and $p_g$
for its population share. A treated-cell average weights every supported
post-adoption cell by its cohort share. A group average first averages within
each cohort over selected post periods and then weights cohorts. A dynamic
average first aggregates cohorts at each relative time and then averages the
selected nonnegative relative times. A calendar average first aggregates
cohorts at each calendar time and then averages those calendar effects.
The inspected implementation contains separate branches for all four targets;
its public aggregation function defaults to the group target. Where weights
are estimated, their influence contributions accompany the effect estimates.
\citep{exec24_panelr_did}

In an authored panel, cohort one has effects two and four at times one and
two; cohort two has effect ten at time two. Their population shares are
$3/4$ and $1/4$. The simple treated-cell target is
\[
 \frac{(3/4)2+(3/4)4+(1/4)10}{3/4+3/4+1/4}=4.
\]
The group target is $(3/4)(2+4)/2+(1/4)10=19/4$.
Dynamic effects at horizons zero and one are both four, giving an overall
value four. Calendar effects are two and $11/2$, averaging to $15/4$.
At a stipulated 100 USD per outcome unit, these planning quantities are 400,
475, 400 and 375 USD per respective target unit. The differing denominators
must be explicit before a sponsor treats them as comparable returns.
\citep{exec24_panelr_contract}

Balancing on one post-adoption lag retains only the early cohort in this
example, giving the path $(2,4)$ and mean three. That is a supported,
restricted population, not an estimate of the late cohort's absent longer
history. In the selected source, event balancing affects dynamic aggregation;
simple and group aggregation warn that they ignore that option. Calendar
aggregation warns that it ignores the event-window and balancing options.
The CLI/MCP result must report the resolved request and target population,
not merely repeat options the caller supplied.
\citep{exec24_panelr_did,exec24_panelr_contract}

\subsection{Missing cells and changing weights are population choices}
The selected cell wrapper catches a two-by-two estimation error, warns, and
leaves that cell unavailable. Aggregation rejects missing effects by default;
when removal is requested it filters cells and, in the group branch, removes
groups without retained post-treatment effects in the requested window. The
result therefore describes what remains. It must not be relabeled the original
portfolio without an argument for the missing contributions.
\citep{exec24_panelr_did}

If the late cohort's time-two cell is unavailable in the authored example,
the simple and group averages among remaining cells are both three. The
stipulated complete-data targets were four and $19/4$. Those complete values
are known only because this is a constructed example; a production analysis
must not impute the missing effect by copying them. Requested, eligible,
attempted, estimated and omitted cohort-time cells belong in separate counts.
\citep{exec24_panelr_contract}

The selected \texttt{att\_gt} documentation and cell code also distinguish
sampling-weight clocks. A balanced panel's default two-period comparison
uses earlier-period weights; explicitly varying weights switches to a stacked
repeated-cross-section route while holding covariates at the earlier period.
Base-period or first-period anchoring on an unbalanced panel can remove units
absent at that anchor. For effects $(0,10)$, weights $(9,1)$ give one while
weights $(1,9)$ give nine. If only the first unit is present at an anchor,
the retained target is zero. Neither a constant variable name nor a successful
fit establishes that the population stayed fixed.
\citep{exec24_panelr_did,exec24_panelr_contract}

The stacked route folds pre- and post-period influence contributions back to
the same unit with a factor of one half. Authored stacked contributions
$(-2,2,-4,4)$ become $(-3,3)$; summing without the half gives $(-6,6)$,
four times the variance and twice the standard error under the same variance
calculation. The selected source includes the half. This is a positive
semantic check and a proposed conversion oracle, not a reported defect.
\citep{exec24_panelr_did,exec24_panelr_contract}

\subsection{Controls, anticipation and uncertainty are resolved outputs}
The preprocessing path validates unit-time uniqueness and constant first
adoption within a unit, and treats the intervention as irreversible. A
first-adoption field cannot faithfully encode repeated switches between
policies. The not-yet-treated comparison requires a cohort to adopt strictly
after the comparison period plus anticipation. With time two and one period
of anticipation, a cohort adopting at three is not eligible; a cohort adopting
at four is. When no never-treated group exists, the selected preprocessing
can use the last cohort as comparison and removes unsupported later periods.
These transformations and support losses must appear in the result.
\citep{exec24_panelr_did,exec24_panelr_contract}

Disabling influence-function computation returns point estimates only and
prevents the documented aggregation and inference paths. With inference
retained, the selected multiplier bootstrap permits the individual sampling
unit and at most one additional coarser clustering variable. Its clustered
path aggregates influence contributions to cluster sums, with corresponding
normalization. It also removes degenerate dimensions and nonfinite bootstrap
maxima before forming its simultaneous critical value. A future worker must
return the actual dimensions and denominator used, and qualify sparse or
unstable cases rather than calling every finite critical value calibrated.
\citep{exec24_panelr_did}

An aggregation request cannot introduce a cluster variable absent from the
stored group-time object. The inspected code warns and removes unsupported
extra clustering. Thus a request for branch clustering on an object retaining
only region clustering cannot be reported as fulfilled. The adapter should
refuse or explicitly mark the result nonconformant and request recomputation
from the appropriate first-stage object. A simultaneous-band request has a
similar contract: invalid or unexpectedly small bootstrap critical values
trigger a pointwise fallback, and the returned status must say pointwise.
A large critical value triggers a warning rather than proving useful precision.
\citep{exec24_panelr_did,exec24_panelr_contract}

\subsection{Choose among complementary R routes with their assumptions intact}
The previously inspected \texttt{DRDID} dispatcher separates panel and
repeated-cross-section designs and improved from traditional nuisance routes.
Its selected multiplier helper does not itself accept arbitrary cluster
structure. \texttt{did} can organize such two-period comparisons, aggregation
and a separate influence-based covariance layer, but that composition needs
a qualified dependency version and identity alignment. The \texttt{did2s}
route fits an untreated first stage and a full-sample second stage, with
its own correction for generated outcomes and clustered inference. Ordinary
second-stage regression uncertainty is not automatically that correction.
\citep{exec24_drdid_reuse_source,exec24_did2s_reuse_source}

\texttt{fixest} supplies a broad projection and covariance engine. Its
selected defaults, singleton/perfect-fit removals, small-sample scaling,
covariance repair and event-cohort mapping require explicit decisions.
Fixed effects alone do not select clustered errors. Its event-study
aggregation uses full covariance, which a summary-table adapter must preserve.
\texttt{HonestDiD} adds sensitivity to specified trend deviations, with
relative-magnitude and smoothness restrictions of different meanings. It
requires correctly ordered coefficients and covariance and honest numerical
status; it cannot verify that the imposed restriction is substantively true.
\citep{exec20_fixest_source,exec22_honestdid_engine}

The original \texttt{synthdid} source reviewed here targets a simultaneous
adoption block. It should not inherit arbitrary staggered-adoption support
from another implementation's name. The selected \texttt{gsynth} compatibility
wrapper delegates to \texttt{fect}; the earlier inspection found that its
weight argument was not forwarded even though its documentation described a
weighted effect. The inspected \texttt{fect} engine distinguishes fitting
weights, aggregation weights, factor assumptions, reversal and inference
gates. A stable numerical factor fit does not establish counterfactual
stability or solve an unidentified missing trajectory.
\citep{exec23_synthdid_reuse_source,exec21_gsynth_wrapper,exec21_gsynth_engine}

For a branch rollout, the comparison sequence is therefore to specify target,
clock, treatment history, controls, support, sampling population and permissible
trend restrictions; select a concrete route; retain actual samples and
uncertainty; and compare against independent public or generated oracles.
The qualified scope may be a restricted cohort or a refusal. A software census
also needs explicit tails: alternative panel estimators, reversal designs,
spatial/network dependence and other R packages remain to be screened. The
selected seven families do not establish exhaustive coverage.
\citep{exec24_panelr_contract}

\subsection{Public qualification and costs before expensive bank experiments}
Known-law panels can compare cohort, event and calendar targets; anticipation,
missing cells, sampling-weight anchors, influence-vector ordering, cluster
membership and fallback behavior. A preregistered benchmark must preserve
requested and delivered targets, complete failures, dependence assumptions,
coverage uncertainty and compute costs. Illustrative total cost of 3,000 USD
across 150 requests, 125 conformant outputs and 100 useful outputs gives 20,
24 and 30 USD per denominator. These are planning values, not observed prices.
\citep{exec24_panelr_contract}

The scientific owner approves causal and sampling restrictions; the adapter
owner preserves transformations, objects and refusal semantics; the
platform/rights owner pins the installed dependencies and distribution rights;
and the evaluation owner supplies independent oracles and complete accounting.
No selected source inspection alone grants a native integration. Public
experiments can establish merit before bank experiments, but cannot manufacture
untreated trends, missing counterfactuals or universal validity under arbitrary
dependence. Coverage remains incomplete while search tails and those native
qualification contracts remain open.
\citep{exec24_panelr_contract}
